\section{Introduction}

Training language models on their own generated data across multiple iterations leads to \emph{model collapse}---a progressive degradation in both output quality and lexical diversity \citep{shumailov2024ai, alemohammad2024self}. Each generation of synthetic data carries distributional artifacts from its parent model; when these artifacts compound across iterations, the distribution's tails erode and the model converges toward a narrow, repetitive mode \citep{seddik2024how, briesch2023large, guo2024curious}. As synthetic data becomes an increasingly prevalent component of training pipelines, understanding this feedback loop---and identifying strategies that avoid it---has become a pressing concern \citep{dohmatob2025strong, gerstgrasser2024model}.

Diversity-based data curation is a natural mitigation strategy. If model collapse arises from training on redundant or low-quality synthetic samples, then selecting a maximally diverse subset should preserve the distributional richness that iterative training erodes. Submodular maximization over embedding spaces provides a principled framework for this goal: submodular diversity objectives identify representative subsets that maximize geometric coverage of the data distribution \citep{wei2015submodularity, mirzasoleiman2020coresets}. Embedding-based selection has demonstrated strong results in data pruning \citep{sorscher2022beyond} and data-efficient training more broadly \citep{xie2023data, xia2024less, engstrom2024dsdm}. When the model's own representations define the embedding space, diversity selection should---in principle---directly counteract the homogenization that drives collapse.

We test this intuition in a controlled iterative self-training experiment and find the opposite. Representation-Guided Diversity Selection (RGDS)---which selects the most diverse 50\% of synthetic data according to the training model's own last-hidden-layer embeddings---increases Generation~7 perplexity by +72\% relative to random selection under the same data budget ($74.29 \pm 7.71$ vs.\ $43.18 \pm 1.43$, $n{=}3$) and reduces rare n-gram survival by 38\%. The method designed to preserve diversity accelerates collapse instead.

Existing data selection methods---SemDeDup \citep{abbas2023semdedup}, D4 \citep{tirumala2023d4}, DSIR \citep{xie2023data}---are designed for one-shot curation of static datasets where the embedding model is fixed. Active learning accounts for model updates between selection rounds \citep{ash2020deep, kirsch2019batchbald}, but not for the compound effect of selection bias feeding back into the selecting model's own training data across generations. The concept drift literature establishes that representations degrade under distribution shift \citep{lu2019learning, gama2014survey}; however, the interaction between this degradation and embedding-based selection in a self-consuming training loop---where the selecting model is also shaped by its selections---has not been quantified.

We trace this failure to the \emph{embedding alignment gap}: a progressive divergence between what the model's internal representations treat as ``diverse'' and what constitutes genuine content diversity. We quantify this divergence with the \emph{selection ratio}---the mean rare n-gram density of the selected subset divided by that of the rejected subset. A ratio above~1 indicates that the selected samples are richer in rare n-grams than the rejected samples; below~1, the selection preferentially retains content-poor samples. For RGDS, the selection ratio declines from ${\sim}1.0$ at Generation~0 to $0.64 \pm 0.20$ at Generation~7 ($n{=}3$): by the final generation, the model's internal notion of diversity has inverted relative to content diversity.

The sharpest evidence for the embedding alignment gap comes from comparing methods that share identical selection algorithms but differ in their embedding source. All methods that bypass internal embeddings cluster in a narrow band at Generation~7: Random~50\% ($43.18 \pm 1.43$), External-only using frozen Contriever \citep{izacard2022unsupervised} embeddings ($46.36 \pm 2.97$), and Uncurated ($47.04 \pm 2.90$)---all $n{=}3$ (ANOVA $p = 0.219$). Only RGDS produces catastrophic outcomes ($74.29 \pm 7.71$). The External-only method is particularly informative: it applies the same diversity selection algorithm and 50\% budget as RGDS but uses a frozen encoder whose embeddings do not co-degrade with training. Its normal-range performance confirms that the selection algorithm itself is not at fault---the failure arises specifically from using co-degrading internal embeddings as the selection criterion, a form of Goodhart's law in representation space \citep{gao2023scaling, karwowski2024goodhart}.

Our contributions are:
\begin{enumerate}
    \item We quantify the \emph{embedding alignment gap} via the \emph{selection ratio} and show that it worsens progressively, declining from ${\sim}1.0$ to $0.64 \pm 0.20$ over 8 generations of internal-embedding selection ($n{=}3$).
    \item We isolate the failure to internal embeddings specifically: RGDS is the only method producing catastrophic degradation, while random selection, frozen external embeddings, and no selection all cluster in a narrow performance band (Generation~7 PPL 43--50, ANOVA $p = 0.219$, $n{=}3$).
    \item We establish that the alignment gap causes accelerated collapse: under the same 50\% data budget, RGDS increases perplexity by +72\% and reduces rare n-gram survival by 38\% relative to random selection ($n{=}3$), with preliminary evidence of the same pattern at 355M parameters.
\end{enumerate}